Source:
Analysis III (MA 321)
Lecture notes and course material.

Instructor: Prof. A.K. Nandakumaran

Area:
Distribution Theory

Type:
Foundational Notes

Status:
Personal study notes.

Date:
February 2026
\section*{Foundations of Distribution Theory}
These notes summarize core definitions and structural results concerning test functions, distributions, support, order, and distributional convergence. 

\section{Test Functions}
\begin{definition}[$\mathcal{D}(\Omega$]
    Let $\Omega\subset \mathbb{R}^n$, open, then:
    \[
    \mathcal{D}(\Omega)= C_c^\infty(\Omega)=\{\phi:\Omega\to \mathbb{R}|\quad\text{$\phi$ is smooth and $\operatorname{supp}\subset\subset\Omega$\} }
    \]
    So, $\mathcal{D}(\Omega)$ is a space of smooth functions or test functions, where $\phi\in C^\infty(\Omega)$ with a compact support. All these test functions are non-zero within the support and vanishes smoothly outside the support.
    \end{definition}
    \begin{definition}[$\mathcal{D}_K(\Omega)$]
    Let $\Omega\subset\mathbb{R}^n$ and $K\subset\subset\Omega$ compact. Then
    \[
    \mathcal{D}_K(\Omega)=\{\phi\in C_c^\infty(\Omega):\operatorname{supp}(\phi)\subset K\}
    \]
    The relationship between $\mathcal{D}(\Omega)$ and $\mathcal{D}_K(\Omega)$ is:
    \[
    \mathcal{D}(\Omega)=\bigcup_{K\subset\subset\Omega}\mathcal{D}_K(\Omega)
    \]
    \end{definition}
\begin{definition}[Seminorms]
    Seminorm is defined on $\mathcal{D}_K(\Omega)$ as
    \[
    p_{K,m}(\phi) = \sum_{|\alpha|\le m}\sup_{x\in K}|\partial^\alpha\phi(x)|
    \]
    \begin{remark}
        $\mathcal{D}_K(\Omega)$ endowed with locally convex topology generated by the family of seminorms $\{p_{K,m}\}_{m\in\mathbb{N}}$
    \end{remark}
\end{definition}
\subsection{Convergence in $\mathcal{D}(\Omega)$}
A sequence $\{\phi_j\}\subset\mathcal{D}(\Omega)$, converges to 0 in $\mathcal{D}(\Omega)$ if there exists a compact set $K\subset\subset\Omega$ such that
\begin{enumerate}
    \item $\operatorname{supp}(\phi_j)\subset K$ for all $j$,
    \item for every multiindex $\alpha$,
    \[ \sup_{x\in K}|\partial^\alpha\phi_j(x)|\to 0 \quad \text{as } j\to\infty. 
    \]
\end{enumerate}
More generally, a sequence $\phi_j$ converges to $\phi$ in $\mathcal{D}(\Omega)$, if $\phi-\phi_j\to 0$ in the above sense. 

\section{Distributions}
\begin{definition}[Distribution via Sequential Continuity]
    Let $\Omega\subset \mathbb{R}^n$, open. A distribution on $\Omega$ is a linear functional
    \[
    T:\mathcal{D}(\Omega)\to\mathbb{R}
    \]
    such that whenever a sequence of test functions $\{\phi_j\}\subset\mathcal{D}(\Omega)$ converges to 0 in $\mathcal{D}(\Omega)$, we have
    \[
    \langle T, \phi_j\rangle\to 0 \quad \text{as } j\to\infty
    \]
    A distribution is a sequentially continuous linear functionals on $\mathcal{D}(\Omega)$. 
\end{definition}
\begin{theorem}[Sequential Continuity Equivalent to Seminorm Estimate]
  A linear functional $T:\mathcal{D}(\Omega)\to \mathbb{R}$ on $\mathcal{D}(\Omega)$ is a distribution if and only if for every compact set $K\subset\subset\Omega$, there exists a constant $C>0$ and $m\in\mathbb{N}$, such that
  \[
  |\langle T,\phi\rangle|\le Cp_{K,m}(\phi)
  \] 
\end{theorem}
Hence, sequential continuity is equivalent to seminorm estimates.
\begin{definition}[Order of a Distribution]
    Let $T\in\mathcal{D}'(\Omega)$ is a distribution of order $\le m$ on a compact set $K\subset\subset\Omega$, if there exists a constant $C>0$ for all multiindex $\alpha$ such that
    \[
    |\langle T, \phi\rangle|\le C\sum_{|\alpha|\le m}\sup_{x\in K}|\partial^\alpha\phi(x)|\quad \forall \phi\in \mathcal{D}_K(\Omega)
    \]
\end{definition}
If such an integer $m$ can be chosen uniformly for all compact set $K$, then the distribution $T$ will have finite order, and if not, then the distribution will have infinite order. Hence, order of a distribution depends on the compact set $K$. The seminorm estimate, for $m=0$, reduces to 
\[
|\langle T,\phi(x)\rangle|\le C\sup_{x\in K}|\phi|
\]
This defines a distribution of order 0. For example, if the distribution is $\delta$, then the seminorm estimate becomes:
\[
|\delta(\phi)|=|\phi(0)|\le\sup_{x\in K}|\phi(x)|
\]
Thus, Dirac delta generates a zero-order distribution. 
\begin{remark}
    Distribution order is defined using seminorm estimates. Physically, order of a distribution says how many derivatives needed to control. Thus, more localized a distribution induced by locally integrable function, less derivatives can control the distribution. 
\end{remark}
\section{Locally Integrable Functions as Distributions}

\begin{definition}[$T_f$]
    Let $\Omega\subset\mathbb{R}^n$, open and let function $f\in L_{\text{loc}}^1(\Omega)$. $K\subset\subset\Omega$ be a compact set. We define the linear functional $T_f:\mathcal{D}(\Omega)\to\mathbb{R}$ by
    \[
    T_f(\phi)
    =
    \langle f,\phi\rangle
    =\int_\Omega f(x)\phi(x)\,dx,  \quad \phi \in \mathcal{D}(\Omega)
    \]
    \end{definition}
\begin{proposition}
    If $f\in L^1_{\mathrm{loc}}(\Omega)$,
    then $T_f\in\mathcal{D}'(\Omega)$.
\end{proposition}
    \begin{proof}
Let $\phi \in \mathcal{D}(\Omega)$ and set 
$K = \operatorname{supp}(\phi) \subset\subset \Omega$. 
Since $f \in L^1_{\mathrm{loc}}(\Omega)$, we have
\[
\int_K |f(x)|\,dx < \infty.
\]
Hence
\[
|T_f(\phi)|
=
\left|\int_K f(x)\phi(x)\,dx\right|
\le
\sup_{x\in K} |\phi(x)|
\int_K |f(x)|\,dx,
\]
which shows that $T_f(\phi)$ is well defined.

Now let $\{\phi_j\}\subset \mathcal{D}(\Omega)$ 
converge to $0$ in $\mathcal{D}(\Omega)$. 
Then there exists a compact set $K$ containing 
$\operatorname{supp}(\phi_j)$ for all $j$, and
\[
\sup_{x\in K} |\phi_j(x)| \to 0.
\]
Therefore,
\[
|T_f(\phi_j)|
\le
\left(\int_K |f(x)|\,dx\right)
\sup_{x\in K} |\phi_j(x)|
\to 0,
\]
which proves sequential continuity.

Linearity follows immediately from linearity of the integral. Let $\phi,\psi\in C_c^\infty(\Omega)$ and $\operatorname{supp}(\phi)\subset K$, $\operatorname{supp}(\psi)\subset K$. Then
    \[
    T_f(a\phi+b\psi)=aT_f(\phi)+bT_f(\psi)
    \]
    For $a, b\in\mathbb{R}$. The above equality holds true because $\langle f, \phi\rangle$ and $\langle f, \psi\rangle$ are well defined by sup bound. Thus, $T_f$ follows linearity. Hence, $T_f$ defines a distribution induced by locally integrable function $f$. 
\end{proof}
\begin{remark}
    The above proof for $T_f$ as a distribution, we only needed $m=0$, which means for constant $C_K>0$:
    \[
    |T_f(\phi)|\le C_K p_{K,0}(\phi)
    \]
    Hence, $T_f$ is of order zero, or more generally, any distribution induced by locally integrable function has order zero. This extends the previous remark that $\delta$ (and more generally, any distribution induced by a locally integrable function) has order zero. 
\end{remark}
